\unnumberedChapter{Simulator User Interface}

The simulator is controlled through an interactive command-line interface. Each command is entered at the command prompt and is executed immediately. A command consists of a command name followed by zero or more arguments. Arguments are not limited to literal values; they may be arbitrary expressions containing constants, variables, register values, memory references, and function calls. This chapter introduces the syntax of the command language. The individual commands are described in detail in the following chapter.

\unnumberedSection{Command Line}

The simulator displays a command prompt before accepting each command. The prompt contains a command counter which uniquely identifies every command entered during the session.

\begin{list}{}{\leftmargin=2em}
\item
\begin{verbatim}
(100) ->
\end{verbatim}
\end{list}

A command consists of a command name followed by an optional list of comma-separated arguments.

\begin{list}{}{\leftmargin=2em}
\item
\begin{verbatim}
command arg1, arg2, arg3
\end{verbatim}
\end{list}

Each argument may be a literal value or an arbitrary expression.

Long commands may be split over several input lines by terminating each continued line with a backslash (\verb|\|). The simulator concatenates the lines before parsing the command.

\begin{list}{}{\leftmargin=2em}
\item
\begin{verbatim}
WRITE R1, \
      R2 + 4, \
      "Hello"
\end{verbatim}
\end{list}

The simulator maintains a history of previously entered commands. The \texttt{HIST} command displays the command history, while the \texttt{DO} and \texttt{REDO} commands allow previously entered commands to be executed again or edited before execution.

\unnumberedSection{Data Types}

The command language supports three basic data types:

\begin{smallGapItemize}
\item Numbers. A number is signed 64-bit quantity.
\item Booleans. A boolean has the value "true" or "false".
\item Strings. A string is character sequence enclosed in apostrophes.
\end{smallGapItemize}

\unnumberedSubsection{Numbers}

All numbers are represented as signed 64-bit integers. Numbers may be entered in decimal, hexadecimal, or binary notation. Hexadecimal numbers use the prefix \verb|0x|, while binary numbers use \verb|0b|.

\begin{list}{}{\leftmargin=2em}
\item
\begin{verbatim}
12345
-17

0xFFFF1000
0xFFFF_1000_0010

0b01101001
0b1010_1111_0000
\end{verbatim}
\end{list}

Underscore characters may be inserted between digits to improve readability. They do not affect the numerical value.

\unnumberedSubsection{Booleans}

Boolean values represent logical truth and may assume one of the two values

\begin{list}{}{\leftmargin=2em}
\item
\begin{verbatim}
true
false
\end{verbatim}
\end{list}

Boolean values are typically produced by comparison operators but may also be entered directly.

\unnumberedSubsection{Strings}

Strings are enclosed in double quotation marks.

\begin{list}{}{\leftmargin=2em}
\item
\begin{verbatim}
"Hello World"
\end{verbatim}
\end{list}

Adjacent string literals are automatically concatenated, allowing long strings to span multiple input lines.

\begin{list}{}{\leftmargin=2em}
\item
\begin{verbatim}
"This is a very long "
"string which is "
"assembled automatically."
\end{verbatim}
\end{list}

\unnumberedSection{Expressions}

Most command arguments are expressions. Expressions are evaluated before the command itself is executed.

\begin{list}{}{\leftmargin=2em}
\item
\begin{tabular}{lll}
\textbf{Operator} & \textbf{Operand Type} & \textbf{Description} \\ \hline
\verb|+|  & Number           & Addition \\
\verb|-|  & Number           & Subtraction \\
\verb|*|  & Number           & Multiplication \\
\verb|/|  & Number           & Signed division \\[2mm]

\verb|==| & Any              & Equal comparison \\
\verb|!=| & Any              & Not equal comparison \\
\verb|<|  & Number           & Less than \\
\verb|<=| & Number           & Less than or equal \\
\verb|>|  & Number           & Greater than \\
\verb|>=| & Number           & Greater than or equal \\[2mm]

\verb|&|  & Number           & Bitwise AND \\
\verb!|!  & Number           & Bitwise OR \\
\verb|^|  & Number           & Bitwise exclusive OR \\[2mm]

\verb|&&| & Boolean          & Logical AND \\
\verb+||+ & Boolean          & Logical OR \\
\verb|!|  & Boolean          & Logical NOT \\
\end{tabular}
\end{list}

Parentheses may be used to explicitly control the order of evaluation. Typical expressions are

\begin{list}{}{\leftmargin=2em}
\item
\begin{verbatim}
R1 + 4
(R2 + R3) * 8
(R5 & 0xff) == 0x80
(R1 + [0x1000]) / 4
\end{verbatim}
\end{list}

\unnumberedSection{Predefined Variables}

The simulator provides a collection of predefined variables describing the current machine state. These include the general-purpose registers, control registers, processor status, and other architectural state.

\begin{longtable}{@{}p{0.2\linewidth}p{0.2\linewidth}p{0.3\linewidth}@{}}
    \caption{Predefined Variables} \\
    \toprule
    \textbf{Name} & \textbf{Type} & \textbf{Purpose} \\
    \midrule
    \endfirsthead
    \toprule
     \textbf{Name} & \textbf{Type} & \textbf{Purpose} \\
    \midrule
    \endhead
    \midrule
    \multicolumn{2}{r}{\textit{Continued on next page}} \\
    \midrule
    \endfoot
    \bottomrule
    \endlastfoot
    \textbf{R0 .. R15} & numeric & General register. \\
    \midrule
    \textbf{C0 .. C15} & numeric & Control register. \\
    \midrule
    \textbf{IA} & numeric & Instruction address register. \\
    \midrule
\end{longtable}

Additional predefined variables may be listed using the \texttt{ENV} command. 

\unnumberedSection{User Defined Variables}

User-defined variables may also be created. Once defined, they may be used wherever expressions are accepted.

\begin{list}{}{\leftmargin=2em}
\item
\begin{verbatim}
(8) -> ENV a_var, 200
\end{verbatim}
\end{list}

\unnumberedSection{Predefined Functions}

The command language provides a number of predefined functions. Functions may appear anywhere within an expression and return values which may be used by the enclosing command.

For example,

\begin{list}{}{\leftmargin=2em}
\item
\begin{verbatim}
(9) ->w ASM ("Add r1, r2, r3" )
 0x4410600
 
(10) ->w ADD_OFS( 0x2000_0000, 0x1000 )
 0x20001000
\end{verbatim}
\end{list}

The complete list of predefined functions is given in the command reference.

\unnumberedSection{Memory References}

Memory contents may be accessed directly within expressions. The general syntax is

\begin{list}{}{\leftmargin=2em}
\item
\begin{verbatim}
[ [type] address ]
\end{verbatim}
\end{list}

where \textit{address} is any expression evaluating to a memory address and \textit{type} specifies the access size.

The supported access types are

\begin{list}{}{\leftmargin=2em}
\item
\begin{tabular}{ll}
BYTE   & 8-bit access \\
HALF   & 16-bit access \\
WORD   & 32-bit access \\
DOUBLE & 64-bit access
\end{tabular}
\end{list}

If no type is specified, a doubleword access is assumed.

Examples:

\begin{list}{}{\leftmargin=2em}
\item
\begin{verbatim}
[0x1000]
[BYTE 0x1000]
[WORD R1 + 8]
[DOUBLE StackPointer]
\end{verbatim}
\end{list}

Memory references may appear wherever expressions are accepted.

\unnumberedSection{Example Session}

// ??? needed ? what to show ?

\begin{list}{}{\leftmargin=1em}
\item
\begin{verbatim}
Twin-64 Simulator, Version: A.01.00, Patch Level: 122
 Git Branch: main
 Load Config File: "./Twin64/config/T64ConfigFile"
 Open Log File: "./Twin64/log/T64LogFile"
 (0) ->
 (1) -> 

\end{verbatim}
\end{list}	
